\documentclass[10pt,a4paper]{article}
\usepackage[margin=21mm,headheight=15pt]{geometry}
\usepackage[T1]{fontenc}
\usepackage{lmodern}
\usepackage{microtype}
\usepackage{longtable,array,booktabs}
\usepackage{enumitem}
\usepackage{xcolor}
\usepackage{fancyhdr}
\usepackage{needspace}
\usepackage[colorlinks=true,linkcolor=blue!45!black,urlcolor=blue!45!black]{hyperref}
\definecolor{ink}{HTML}{193F50}
\setlength{\parindent}{0pt}
\setlength{\parskip}{5pt}
\setlength{\emergencystretch}{2em}
\renewcommand{\arraystretch}{1.17}
\setlist{itemsep=4pt,topsep=4pt,leftmargin=1.7em}
\pagestyle{fancy}
\fancyhf{}
\fancyhead[L]{\small\color{ink}SealBot / CareBot}
\fancyhead[R]{\small Design revision 2026-09-05}
\fancyfoot[C]{\small\thepage}
\begin{document}
\begin{titlepage}
\color{ink}
\vspace*{20mm}
{\large ROBOTICS FOR GOOD 2026--2027\par}
\vspace{9mm}
{\Huge\bfseries Two robots. Nine days.\par}
\vspace{7mm}
{\Large SealBot and CareBot\par}
\vspace{12mm}
\rule{\linewidth}{1pt}
\vspace{8mm}
{\large Rs 6,700 planned / Rs 7,000 maximum\par}
{\large ESP32-CAM and simple mechanisms for a school team\par}
\vspace{12mm}
\textbf{SealBot}\par
Collect a floor-staged beam $\rightarrow$ guide it into position $\rightarrow$ release
$\rightarrow$ repeat for the second beam $\rightarrow$ park.\par
\vspace{6mm}
\textbf{CareBot}\par
Deliver all kits $\rightarrow$ capture one patient $\rightarrow$ inspect under fixed light
$\rightarrow$ deliver by color $\rightarrow$ repeat while time remains.\par
\vfill
\color{black}
Budget revision and nine-day build plan\par
5 September 2026\par
Based on the supplied Rulebook-2026.pdf and official camera documentation.\par
Borrowed tools and practice surface assumed. Sample handling is excluded.\par
Prototype dimensions and timing targets require physical validation.
\end{titlepage}
\tableofcontents
\clearpage
\needspace{12\baselineskip}\section{The decision}
Yes, ESP32-CAM can work. Use it to classify one stopped patient in a shaded pocket, with fixed white lighting and local image processing. A Nano handles the motors and servos. This is the camera task I recommend for this budget and deadline.

The previous Pixy2, carousel, elevator, and sample indexer design was too ambitious for two robots built by 11th graders in nine days. Replacing only the camera would not fix that. This revision removes those mechanisms and targets reliable partial completion instead of promising every mission.

Build SealBot to place the two beams. Build CareBot to deliver all ten kits and then transport one patient at a time. Initially target two to four correct patient deliveries; increase that only after timed trials. Sample collection and laboratory insertion are outside this nine-day build. That sacrifices potential points and may incur sample penalties, but protects the simpler scoring tasks.

Budget Rs 6,700 for the plan below, with another Rs 300 left below the Rs 7,000 ceiling. This assumes a borrowed laptop, soldering iron, multimeter, drill, hand tools, and flat practice surface. It assumes four to six students can work about four focused hours daily, with school supervision for tools and batteries. The electronics must be available by Day 2. These are planning assumptions, not confirmed resources.

No unbuilt design can honestly be guaranteed fault-free. The main remaining risks are beam release geometry, navigation to small patients, camera calibration, and late parts. The day-by-day acceptance gates below make those risks visible early.

\needspace{12\baselineskip}\section{What we keep and what we remove}
\small\begin{longtable}{@{}>{\raggedright\arraybackslash}p{0.22\linewidth}>{\raggedright\arraybackslash}p{0.42\linewidth}>{\raggedright\arraybackslash}p{0.27\linewidth}@{}}
\toprule
\textbf{Previous design fault} & \textbf{Replacement} & \textbf{Why it helps} \\
\midrule\endfirsthead
\toprule\textbf{Previous design fault} & \textbf{Replacement} & \textbf{Why it helps} \\\midrule\endhead
Pixy2 consumes more than the total budget & One ESP32-CAM with OV2640 and verified PSRAM & Enough processing for a fixed inspection area, at a small fraction of the cost. \\[3pt]
Twelve-pocket carousel needs indexing and transfer interlocks & One patient pocket & No inventory map, storage indexing, or mixed-color bin. \\[3pt]
Elevator creates two extra handoffs & Floor-level shallow scoop & Patient stays in one place from pickup to delivery. \\[3pt]
Three sample trays and precision laboratory insertion & Omit sample mission for this build & Removes an entire mechanism and its 2 mm nominal radial alignment problem. \\[3pt]
Stacked or diagonal beam carriage creates fit and unloading problems & Beams start on the floor; a low clamp moves one at a time & No carrier, lifting, gravity slide, or support withdrawal. \\[3pt]
Metered kit magazines need holdback and counting & Three pre-counted batch chutes & Each destination receives its whole assigned load in one operation. \\[3pt]
Timed driving drifts with battery and payload & Cheap wheel pulse sensors plus field references & Measures wheel travel and corrects at existing lines and walls. \\[3pt]
Fixed robot delay assumes nothing jams & Separate working areas, bumpers, and timeout stops & A delayed robot cannot authorize blind movement by the other. \\[3pt]
\bottomrule\end{longtable}\normalsize
There are six working servos: one beam clamp, three kit gates, one patient gate, and one short patient ejector. Buy two spares. There are four drive motors, two Nanos, and only one camera. There are no lifting motors, turntables, custom PCBs, or required 3D prints.

\needspace{12\baselineskip}\section{Budget for both robots}
These are spending allowances including expected GST, not a checked-out supplier cart. The price references in Section 11 explain the main allowances. Small hardware, local materials, freight, and stock substitutions remain estimates. Recalculate the landed total before ordering; do not buy the cheapest version if it lacks a camera, PSRAM, motor pair, or battery protection.

\small\begin{longtable}{@{}>{\raggedright\arraybackslash}p{0.22\linewidth}>{\raggedright\arraybackslash}p{0.42\linewidth}>{\raggedright\arraybackslash}p{0.27\linewidth}@{}}
\toprule
\textbf{Item} & \textbf{Quantity / allowance basis} & \textbf{Total Rs} \\
\midrule\endfirsthead
\toprule\textbf{Item} & \textbf{Quantity / allowance basis} & \textbf{Total Rs} \\\midrule\endhead
2WD kits, motors, wheels, caster, brackets & 2 at 350 & 700 \\[3pt]
Nano-compatible CH340 boards & 2 at 200 & 400 \\[3pt]
ESP32-CAM with OV2640 & 1, including tax allowance & 720 \\[3pt]
USB-to-TTL programmer and programming leads & 1, 3.3 V UART logic & 150 \\[3pt]
DRV8833 dual motor-driver modules & 2 at 80 & 160 \\[3pt]
SG90 180-degree servos & 8 at 100; six working, two spare & 800 \\[3pt]
Five-channel line sensor boards & 2 at 190; initially use three channels & 380 \\[3pt]
Wheel slot sensors and encoder discs & 4 sensor/disc sets & 200 \\[3pt]
Preassembled protected 2S Li-ion packs & 2 at 550, about 2 Ah each & 1,100 \\[3pt]
Compatible 8.4 V CC/CV charger & 1, charge packs sequentially & 600 \\[3pt]
LM2596-type adjustable regulators & 4 at 60 & 240 \\[3pt]
Switches, fuses, wire, connectors, perfboard, LEDs, capacitors, resistors & Shared allowance & 350 \\[3pt]
Deck offcuts, PET sheet, hinges, fasteners, tape, glue, guards & Shared allowance & 250 \\[3pt]
Practice game pieces and markings & School-cut offcuts; surface borrowed & 150 \\[3pt]
Delivery / local collection & Consolidated orders & 300 \\[3pt]
Small repair reserve & In addition to spare servos & 200 \\[3pt]
Planned total & Both robots and shared supplies & 6,700 \\[3pt]
Unallocated headroom & Keep below the hard ceiling & 300 \\[3pt]
Maximum authorized design budget & Do not spend beyond this & 7,000 \\[3pt]
\bottomrule\end{longtable}\normalsize
The Rs 5,000 end of the range requires about Rs 1,700 of reuse or borrowing. Two existing protected packs and a compatible charger would account for that allowance. I would not claim a new-from-scratch Rs 5,000 build by excluding charging, wiring, or delivery costs.

Use school stock or one local supplier where possible. A Rs 20 saving is not useful if it adds another delivery fee or three days of waiting. If prices exceed the plan, first reuse chassis material, programmer, or power equipment. Do not remove fuses or substitute an unprotected loose-cell pack to meet the number.

\needspace{12\baselineskip}\section{CareBot: kit delivery and one patient at a time}
\needspace{10\baselineskip}\subsection{Three kit chutes}
Preload three covered chutes with two, six, and two kits. Each has one servo-operated gate. Open the matching gate only when the outlet is safely inside its destination. Since the entire batch leaves together, there is no single-kit escapement or software counting mechanism.

Build channels from smooth PET or thin plywood lined with smooth plastic. Start with roughly 28-30 mm internal clearance for a 25 mm kit and an adjustable slope. Increase the slope only until all pieces slide reliably. Support the load on the chute and a mechanical gate stop; the servo should release the latch rather than carry the kit weight continuously.

The six-kit hospital channel may be longer than a PCC channel. Mock up its complete loaded volume before mounting the camera. Avoid a wide hopper where two cubes can bridge across the outlet. Keep each chute single-file, covered, and free of inward-facing screw heads.

Deliver each destination's full batch on one visit. Test the outlet offset and retreat so no kit remains on the boundary or gets dragged out. An open gate is not evidence of a complete delivery; inspect all ten pieces after every practice run. If one chute sticks, fix it on the bench rather than adding repeated shaking to the route.

\needspace{10\baselineskip}\subsection{A small scoop with positive release}
Use a shallow front scoop with a 45-50 mm gathering mouth narrowing to an approximately 24-26 mm pocket. Start with a rounded 0.3-0.5 mm PET lip and a low floor. Adjust against the actual painted 20 mm cylinders. Side walls keep the captured patient in place without squeezing it.

One SG90 closes the front gate after the cylinder seats. A second SG90 drives a short padded ejector. This extra inexpensive actuator is justified: opening a gate alone does not guarantee that a patient slides off the scoop. At the correct destination, open the gate, push the patient fully out, retract the ejector, and reverse clear.

A very light lever switch in the pocket confirms seating; it must not require enough force to topple the cylinder. The camera also checks that the expected inspection region contains usable colored pixels. A partly seated or tipped cylinder gets one gentle reseating attempt, then UNKNOWN. Do not continue scooping into a cluster while a patient is onboard.

The hood, gate, ejector, camera, and scoop remain fixed to the robot. No parts are deposited on the field. Test the lip for scratching and the ejector for knocking patients outside the zone.

\needspace{10\baselineskip}\subsection{Delivery strategy}
After kits, approach one accessible patient, capture it, stop and classify it, then drive directly to its destination. Red goes to H, yellow to a PCC, and green to RZ. Prefer the nearer PCC for yellow; attempt a two-and-two split only if time and the confirmed deliveries justify another trip.

One-at-a-time transport costs more travel than batch sorting, but eliminates almost all storage failures. Target two to four successful deliveries for the first competition configuration. Do not claim twelve patients fit in the time limit without measuring every round trip.

If a reading stays UNKNOWN after one re-read, retain that patient and park safely. Do not guess a destination or eject it onto the field. This ends triage for that run, but the kit mission remains complete. The possible retained-patient penalty is unresolved in the supplied rules and must be confirmed.

\needspace{12\baselineskip}\section{ESP32-CAM implementation}
\needspace{10\baselineskip}\subsection{Keep the camera job small}
Use an AI-Thinker-style ESP32-CAM with OV2640 and working PSRAM. Verify the delivered board, not just its product title. Mount it rigidly above the scoop, looking at the cylinder's painted top or a broad painted face. Start with an adjustable optical distance around 80-120 mm and establish focus on the fixture. These are trial dimensions, not a lens specification.

Use a matte dark hood and two diffuse white LEDs with resistors. Leave the inspection lighting on at a fixed level. Avoid the harsh onboard flash and direct reflections. The camera should see predominantly the held cylinder in its classification region, not a red kit cross, colored wire, or floor sticker.

Use the official camera driver, JPEG at QVGA as a starting configuration, one frame buffer, and local RGB decoding. A 320 x 240 RGB888 image occupies 230,400 bytes before other buffers, so verify PSRAM allocation and return every camera frame buffer after use. This is an initial configuration to benchmark, not a guaranteed frame rate. Espressif documents memory-bandwidth caveats for RGB/YUV capture, particularly with Wi-Fi. See T1.

Disable Wi-Fi and Bluetooth during the run. No cloud, neural network, live video stream, face recognition, or laptop processing is needed. The Nano asks for a color only while stationary with the pocket occupied.

\needspace{10\baselineskip}\subsection{Classification that students can debug}
\begin{enumerate}
\item Capture labeled images from all twelve actual cylinders, several rotations, an empty pocket, partial pickups, and different room lighting. Keep some captures separate for validation instead of tuning on everything.
\item Crop a fixed inspection region. Convert its pixels to HSV or normalized RGB. Reject pixels that are too dark, too bright, or insufficiently saturated. Handle red on both sides of the hue wraparound.
\item Count red, yellow, and green votes using thresholds calibrated on the fixture. As an initial rule, require at least 60 percent usable pixels in the crop and an 80 percent winning share among usable pixels. These thresholds are starting values; the shares are not probabilities of correctness.
\item Require the same accepted color across five fresh frames acquired after the request. Reject empty or stale images. Allow one repeat after settling, then return UNKNOWN. Target completion within one second; cap the whole inspection attempt at two seconds and verify actual timing on Day 2.
\end{enumerate}
Let automatic exposure settle during calibration, then hold exposure, gain, and white balance steady where the sensor API supports it. Recalibrate under the installed hood if the image is clipped or dark. Do not paste universal hue numbers from a tutorial and assume they fit the competition paint.

\needspace{10\baselineskip}\subsection{Camera-to-Nano connection}
Use a wired UART request and response with a request ID, color, and checksum. Only accept a response for the current request. Boot messages and malformed packets are ignored. A two-second timeout produces UNKNOWN, never a default color.

With microSD disabled, GPIO13 as camera UART RX and GPIO14 as TX are candidate free pins on the standard AI-Thinker layout; confirm the actual board pinout. Nano TX is 5 V and must pass through a divider, for example 1 kOhm series and 2 kOhm to ground, before ESP RX. ESP TX is 3.3 V; verify the Nano's input-high requirement for the board supplied. Join grounds.

Provide disconnectable UART links for programming. Do not connect a programmer and the Nano as competing transmitters. Use 3.3 V UART logic on the programmer and a stable regulated 5 V supply at the camera board's 5 V input. Never apply 5 V directly to an ESP GPIO. Remove the boot strap after flashing.

\needspace{10\baselineskip}\subsection{Finding a patient is a separate problem}
A camera in a pocket does not automatically solve navigation to randomly positioned patients. Start from measured sticker coordinates and scan a small local area with the empty scoop. The approach view may supply a coarse colored-region hint, but final classification occurs only after capture.

Confirm with the organizer whether Senior randomization changes only color assignment or also allowed positions. If broad position changes are legal and the local scan cannot acquire cylinders reliably by Day 5, restrict triage to observed reachable candidates or disable it. Do not let an unproven field-search algorithm consume kit-delivery time.

\needspace{12\baselineskip}\section{SealBot: collect floor-staged beams one at a time}
Page 4 explicitly permits the beams to start on the board inside the starting zone. Use that provision. Stage them upright in two parallel rows, with a gap for the gripper. SealBot collects the nearer beam, places it, returns for the second, then parks. This trades one extra trip for removing the whole carrier mechanism.

Use a low U-shaped clamp across the front. Two padded rear contact points push the beam near its base; one servo closes a linked pair of retaining fingers around the opposite long face. The field supports the beam's weight. The fingers retain it during slow reversing and turning without lifting it. Make the contact height adjustable, initially around 5-10 mm above the field, and avoid any part dragging underneath the wood.

Pushing high on a 60 mm-tall beam standing on a 20 mm base can tip it. Low contacts, restrained acceleration, and gentle turns reduce that risk. Test the actual field surface for sliding resistance and marking. If the beam tips or damages the surface, this mechanism has not passed; do not assume more clamping force will solve it.

At the target, approach so the beam's end reaches its fixed wall, then guide the beam along that wall toward the marked edge. The robot's pusher faces the beam's long side. Finish with the beam touching the line and required corner contact, open the fingers fully, and back away. No support has to be pulled from underneath it. Avoid rotating beside an already placed beam.

Choose placement order in a full-size corner fixture. The second beam must touch the first and its fixed wall. Confirm that the robot remains outside the enclosure and that open fingers clear the joint. Target one repeatable beam before the paired closure. If closure fails, keep the tested single-beam route.

\needspace{10\baselineskip}\subsection{Starting arrangement}
For a packaging mock-up, reserve a 300 x 60 mm strip along one edge of the 480 x 280 mm start area for both upright beams. It accommodates the 280 mm length, two 20 mm bases, and approximately 20 mm between them. Place the two robots in the remaining strip, each targeting at most 220 x 200 mm including mechanisms, with a 10 mm gap. The nominal occupied depth is 60 + 10 + 200 = 270 mm and robot-row width is 450 mm.

Measure all preload, camera brackets, gates, and decoration in that arrangement. The dimensions are targets, not verified CAD. Keep gripper swing clear of the second beam. Orient SealBot toward the beam strip at setup so it can acquire the first beam without an in-place turn between the two robots. CareBot exits away from the beam strip first.

After departure, SealBot's held beam makes its moving envelope wider than the chassis, up to the beam length plus gripper clearance. The route must clear this entire width. If a stock chassis is too long with the scoop or clamp fitted, reuse its drive parts on a shorter offcut deck.

The second pickup uses the starting area, which also contains RZ. Reserve that area for SealBot early in the match. End its second-pickup attempts by 65 seconds, then clear and park outside RZ; CareBot does not plan green delivery before 75 seconds. These windows reduce conflict but do not prove clearance if a robot stalls. CareBot must stop on obstruction instead of forcing its way into the area.

Since samples are omitted, SealBot can close without first visiting the laboratory. Do not add sample collection after closure or plan to drive over beams. Adding samples later would require another sequence and validation cycle.

\needspace{12\baselineskip}\section{Electrical and navigation details that prevent avoidable failures}
\needspace{10\baselineskip}\subsection{Two power branches per robot}
Use one preassembled protected 2S Li-ion pack per robot and the correct 8.4 V lithium CC/CV charger. Confirm protection and charging requirements with the exact pack supplier. A protection board is not a charger, and a TP4056 single-cell module is not a 2S charger. Do not assemble loose lithium cells for this nine-day project.

From the fused battery supply, use one regulator for clean 5 V logic and one for 5 V actuators. Join grounds at the supply return. The camera and Nano use the logic branch; motors and servos use the actuator branch. Move one gate at a time while stopped to reduce peak load. Set both regulators with a multimeter before connecting electronics.

Cheap LM2596 boards do not necessarily deliver their advertised maximum continuously. Measure voltage and temperature on the actual assembled robot. Reject a rail that resets the camera during a motor start. Add local capacitors and motor suppression, shorten power leads, and fix binding mechanisms rather than disabling brownout detection.

Use an accessible red, DC-rated latching power switch per robot that physically disconnects battery power. After restoration, firmware waits for a fresh start-button action. Keep gates mechanically retained without servo torque. Stop all actuation at 118 seconds and never restart the mission automatically after reset.

\needspace{10\baselineskip}\subsection{Drive and sensing}
Use two BO motors per robot, a caster, and one slotted wheel sensor per driven wheel. These cheap single-channel encoders measure pulses, not absolute position or independently verified direction. Infer direction from the drive command and re-reference at field markings. Wheel slip still causes error.

A 65 mm wheel with 20 pulses per revolution gives about 10.2 mm per pulse. That is suitable for approximate travel, not precise beam-joint alignment. Use slow docking and physical references for the last correction. Do not redesign the firmware around millimeter accuracy the sensors do not provide.

A DRV8833 is the planned efficient driver, but package and breakout construction affect its continuous current rating. TI specifies different ratings for different packages. Check measured motor load and brief stall behavior against the actual module. Overheating or repeated protection trips mean the motor load or driver must change. See T2.

Use the existing field tape as references; do not add lines to the competition field. Start with left, center, and right channels of each five-channel board. Use wheel counts between references. A missing expected line triggers a short bounded search, then stop, not indefinite forward motion.

\needspace{30\baselineskip}\subsection{Pin allocation for CareBot's classic Nano}
\small\begin{longtable}{@{}>{\raggedright\arraybackslash}p{0.22\linewidth}>{\raggedright\arraybackslash}p{0.42\linewidth}>{\raggedright\arraybackslash}p{0.27\linewidth}@{}}
\toprule
\textbf{Function} & \textbf{Candidate pins} & \textbf{Check} \\
\midrule\endfirsthead
\toprule\textbf{Function} & \textbf{Candidate pins} & \textbf{Check} \\\midrule\endhead
Camera UART & D0, D1 & Disconnect when programming; level-shift Nano TX. \\[3pt]
Wheel pulse inputs & D2, D3 & Keep interrupt handlers short. \\[3pt]
Motor inputs & D4/D5 and D7/D6 & D5 and D6 provide PWM; implement both directions using the driver truth table. \\[3pt]
Three kit gates & D8, D9, D10 & Servo pulses, not motor PWM. \\[3pt]
Patient gate and ejector & D11, A3 & Five servo signals total. \\[3pt]
Start button & D12 & Require a new press after boot. \\[3pt]
Driver sleep control & D13 & Match the breakout's sleep wiring and keep drive disabled at boot. \\[3pt]
Three line channels & A0, A1, A2 & Treat selected board outputs according to their actual type. \\[3pt]
Left and right soft bumper switches & A4, A5 & Dock slowly; unexpected contact stops travel. \\[3pt]
Pocket presence & A6 & Analog reading with an external resistor; A6 is not a normal digital pin. \\[3pt]
Battery divider & A7 & Divider keeps full pack voltage within ADC limits. \\[3pt]
\bottomrule\end{longtable}\normalsize
The standard AVR Servo library uses Timer1 on this Nano family, so keep motor PWM off D9/D10. Do not change Timer0 frequency on D5/D6 because the match timer depends on normal timekeeping. SealBot uses the same drive map but only one servo output. This is a pin plan, not uploaded or bench-tested firmware. See T3.

\needspace{12\baselineskip}\section{Mission control and collision handling}
Use small state machines driven by elapsed time and sensor events. CareBot runs kits first, then acquire, inspect, deliver, and repeat while delivery time remains. SealBot runs acquire first beam, depart, place, return and acquire second beam, place, clear, and park. Every movement has a timeout and one recovery at most.

Activate both robots at the official start. CareBot departs first; SealBot starts its beam pickup after that exit. Keep SealBot's routes in the lower working area and reserve its early return for the second beam. Test the entire swept envelope of held beams, not just the wheels. Keep its final parking location outside RZ.

Prefer routes that do not cross after departure. Where crossing is unavoidable, neither robot enters an occupied corridor; low-speed bumper contact causes a stop, not a push. Soft bumpers are a last-resort contact stop, not long-range collision sensing. A stalled robot may reduce the score of the other, so run delayed-start and blocked-route tests.

Budget CareBot roughly 35-45 seconds for kits, then use measured round-trip patient times, including any wait for RZ access. If a patient cycle takes 15-20 seconds, only a few fit. Stop acquiring whenever remaining time is less than delivery time from the current position plus a ten-second margin. At 118 seconds stop independently of the camera or current state.

Target SealBot clearing the starting area for good by 65 seconds and finishing by about 70 seconds. These are planning targets, not results. If it cannot collect the second beam before its return window closes, abandon that pickup. No mechanism gets repeated retries for the rest of the match.

\needspace{12\baselineskip}\section{Nine-day schedule}
\begin{enumerate}
\item Day 1: confirm the local rules, existing equipment, and parts availability. Make the starting gauge and beam-release mock-up with actual-size pieces. Order or collect parts immediately. Assign students to mechanics, wiring, and code; keep one shared notebook of dimensions and changes.
\item Day 2: power and flash both Nanos and the camera. Demonstrate red, yellow, green, and UNKNOWN in the fixed pocket. Assemble both drive bases and test power isolation. If parts are still unavailable, reduce the build rather than pretending nine full days remain.
\item Day 3: finish the three batch chutes and the patient gate/ejector. Run 30 stationary releases per chute and 30 patient pickup/release cycles. Correct jams now; do not hide them with longer delays.
\item Day 4: get each robot to repeat basic travel, turns, line reacquisition, and low-speed docking. Complete one beam placement, then work on the second. Verify both loaded robots fit together with cables and decoration.
\item Day 5: complete the kit route and a single patient round trip. Check patient acquisition under the organizer's actual randomization rule. This is the triage decision point: if acquisition or color acceptance is unreliable, stop expanding that feature.
\item Day 6: integrate both robots, set the final route sequence, and add deadline, stall, missing-line, and camera-timeout handling. Freeze mechanisms and parts at the end of the day. No carousel, lift, new camera library, or sample mechanism after this point.
\item Day 7: run at least ten full matches with varied patient colors, lower battery charge, and room-light changes. Log each failure and score. Test one delayed robot and one blocked path. Fix the largest repeated fault first.
\item Day 8: run another ten full matches. Require at least nine correct kit distributions, repeatable unsupported beam placement, no accepted wrong colors, no collisions, and stops before 120 seconds. Disable a failing optional patient routine instead of adding a late workaround.
\item Day 9: rehearse setup, charge packs, label connectors, secure wiring, save known-good firmware, and pack the spare servos and repair materials. Make only calibration adjustments. Do not rebuild the robot the night before the event.
\end{enumerate}
Before enabling triage, test at least 120 labeled pocket presentations across actual cylinders, rotations, partial captures, and empty frames. Target zero wrong accepted classifications and no more than ten percent UNKNOWN. This is a practical screening test, not a statistical guarantee. Keep a separate validation set after tuning.

\needspace{12\baselineskip}\section{Scoring tradeoffs and remaining risks}
The supplied Senior scoring table gives 70 positive points for complete containment and 50 for exact kit distribution. Each correct patient adds five, with extra bonuses for complete color groups. A two-to-four-patient run therefore offers 130-140 positive points before any earned color-group bonus and before penalties. This is an illustrative outcome, not a forecast.

Samples left in quarantine incur minus five each under page 23. Page 21 says minus three for patients outside their correct destination, while page 23 instead says minus five for a patient placed in a wrong destination. Page 7 also says incomplete actions simply do not score. These statements conflict.

Under an interpretation penalizing every undelivered patient at minus three and each untouched sample at minus five, two correct patients with complete beams and kits would give 130 - 30 - 15 = 85 before other bonuses or penalties. Four correct patients would give 140 - 24 - 15 = 101. Confirm the organizer's interpretation before predicting a net score. Do not market this as a 140-point guaranteed robot.

The architecture's remaining high-risk item is paired beam placement. A mock-up must prove low-force sliding, upright transport, clean unclamping, and the corner joint. The camera risk is manageable because the inspected object is stationary; the harder remaining triage problem is reaching it. The electrical risk is supply sag, addressed by separate rails and loaded tests. The schedule risk is delivery, addressed by local sourcing and a Day 6 freeze.

If beam closure fails testing, use one reliable beam. If triage fails, keep kit delivery. If both robots interfere, change route ownership or keep the weaker robot parked after its start state. These are explicit fallback configurations, not claims that every subsystem will succeed.

\needspace{12\baselineskip}\section{Price references and technical sources}
Prices below were inspected on 5 September 2026. A displayed price is not a guarantee of stock, compatible contents, or arrival before Day 2. The BOM rounds prices and includes allowances; it is not a set of purchases.

\begin{itemize}
\item P1. \href{https://probots.co.in/2wd-clear-acrylic-smart-robot-chassis-car-kit.html}{Probots 2WD kit}: Rs 349 including GST, with motors and caster. Verify encoder discs or supply them separately.
\item P2. \href{https://robu.in/product-category/microcontroller-development-board/arduino-boards/arduino-board/}{Robu Arduino-compatible board listings}: retrieved listings showed Nano variants around Rs 195-199 including GST. Confirm the selected board and cable.
\item P3. \href{https://directronics.in/product/esp32-cam-wifi-module-camera-module}{Directronics ESP32-CAM}: Rs 610 before 18 percent GST, or Rs 719.80. Confirm OV2640 and PSRAM; programmer is separate.
\item P4. \href{https://robocraze.com/products/sg90-micro-servo-motor}{Robocraze SG90}: Rs 91 including GST when inspected. The site announced shipping resumes on 7 September for orders after its cutoff, so local stock may be more useful for this build.
\item P5. \href{https://www.tomsonelectronics.com/products/5-channel-line-tracking-sensor-module-bfd-1000}{Tomson BFD-1000}: Rs 189 including tax.
\item P6. \href{https://hubtronics.in/lm393-speed-measuring-sensor-module}{Hubtronics slot sensor}: Rs 40 including GST. Allow separately for discs and mounting.
\item P7. \href{https://robodo.in/products/li-ion-battery-7-4v-2200mah-2c-wth-inbuilt-charger-protection}{Robodo protected 7.4 V pack}: displayed Rs 427 including GST, but availability text is ambiguous. Budget Rs 550 per protected replacement and confirm specifications and stock locally.
\item P8. \href{https://goldstarstore.in/products/8-4v-1a-lithium-battery-charger-made-in-india-laptrust}{Goldstar 8.4 V charger}: displayed Rs 550 and a sold-out label. This supports a price allowance only; obtain a compatible available charger or borrow one. Rs 600 is reserved.
\item P9. \href{https://www.ktron.in/product/lm2596-step-down-module/}{KTRON LM2596 module}: Rs 46.80 including GST. The plan allows Rs 60 each and requires load testing.
\item P10. \href{https://robu.in/product-category/dc-motors/motor-drivers/stepper-motor-driver/}{Robu driver listings}: retrieved listing showed DRV8833 at Rs 72 including GST. Verify package and load suitability rather than relying on the headline current rating.
\item T1. \href{https://github.com/espressif/esp32-camera}{Espressif camera driver}: official memory, capture-format, and frame-buffer guidance.
\item T2. \href{https://www.ti.com/lit/ds/symlink/drv8833.pdf}{TI DRV8833 datasheet}: official drive logic and package-dependent current ratings.
\item T3. \href{https://github.com/arduino-libraries/Servo/blob/master/src/avr/ServoTimers.h}{Arduino Servo timer allocation}: official AVR timer selection.
\item R1. Supplied Rulebook-2026.pdf, pages 4-7 for starting conditions and missions, pages 9-16 for objects and geometry, page 17 for robot requirements, and pages 21-24 for Senior scoring and conflicting penalties.
\end{itemize}
Build the fixed-light ESP camera fixture and one reliable mechanical transfer first. Spend the last three days proving the chosen missions, not adding more of them.

\end{document}
